
\clearpage
\section{Elementary Realizations of the Standard Model}
\begingroup
\footnotesize
\setlength{\tabcolsep}{1.9pt}
\begin{longtable}{L{2.03cm}L{2.03cm}L{1.87cm}L{2.34cm}L{6.55cm}}
\toprule
\rowcolor{HMTNavy}\HMTtablehead{Particle} & \HMTtablehead{Charge} & \HMTtablehead{Spin} & \HMTtablehead{Statistics} & \HMTtablehead{Realization HMT--MD} \\
\midrule
\endfirsthead
\multicolumn{5}{l}{\sffamily\scriptsize\color{HMTGray}Continuation of the table}\\[-1pt]
\toprule
\rowcolor{HMTNavy}\HMTtablehead{Partícula} & \HMTtablehead{Carga} & \HMTtablehead{Espín} & \HMTtablehead{Estadística} & \HMTtablehead{Realización HMT--MD} \\
\midrule
\endhead
\midrule
\endfoot
\bottomrule
\endlastfoot
\(e^-\) & \(-1\) & \(1/\allowbreak{}2\) & fermiónica & punto fijo \((5,5)\), vacancia orientada \\
\rowcolor{HMTLight}\(\mu^-\) & \(-1\) & \(1/\allowbreak{}2\) & fermiónica & bloque leptónico de rango 8 \\
\(\tau^-\) & \(-1\) & \(1/\allowbreak{}2\) & fermiónica & bloque leptónico de rango 16 \\
\rowcolor{HMTLight}\(\nu_e\) & \(0\) & \(1/\allowbreak{}2\) & fermiónica & vector de interacción en soporte neutro 20 \\
\(\nu_\mu\) & \(0\) & \(1/\allowbreak{}2\) & fermiónica & vector de interacción en soporte neutro 20 \\
\rowcolor{HMTLight}\(\nu_\tau\) & \(0\) & \(1/\allowbreak{}2\) & fermiónica & vector de interacción en soporte neutro 20 \\
\(u\) & \(+2/\allowbreak{}3\) & \(1/\allowbreak{}2\) & fermiónica & rama up, generación 1, triplete de color \\
\rowcolor{HMTLight}\(d\) & \(-1/\allowbreak{}3\) & \(1/\allowbreak{}2\) & fermiónica & rama down, generación 1, triplete de color \\
\(c\) & \(+2/\allowbreak{}3\) & \(1/\allowbreak{}2\) & fermiónica & rama up, generación 2, triplete de color \\
\rowcolor{HMTLight}\(s\) & \(-1/\allowbreak{}3\) & \(1/\allowbreak{}2\) & fermiónica & rama down, generación 2, triplete de color \\
\(t\) & \(+2/\allowbreak{}3\) & \(1/\allowbreak{}2\) & fermiónica & rama up, generación 3, triplete de color \\
\rowcolor{HMTLight}\(b\) & \(-1/\allowbreak{}3\) & \(1/\allowbreak{}2\) & fermiónica & rama down, generación 3, triplete de color \\
\(\gamma\) & \(0\) & \(1\) & bosónica & sección nula, dos polarizaciones \\
\rowcolor{HMTLight}\(g^a\) & \(0\) & \(1\) & bosónica & adjunta de color, ocho generadores \\
\(W^\pm\) & \(\pm1\) & \(1\) & bosónica & bloque vectorial cargado conjugado \\
\rowcolor{HMTLight}\(Z^0\) & \(0\) & \(1\) & bosónica & bloque vectorial neutral \\
\(H\) & \(0\) & \(0\) & bosónica & sección escalar persistente \\
\end{longtable}
\endgroup
Las antipartículas se obtienen por \(C_M\). Los antiquarks pertenecen al triplete conjugado; \(W^-\) es el conjugado de \(W^+\); fotón, \(Z\), Higgs y gluones poseen las conjugaciones propias de sus representaciones.

El soporte B1 exacto de las clases elementales tiene rangos

\[
 1,\quad8,\quad16,\quad
 3\times20,\quad
 3\times16,\quad
 3\times20.
\]
Therefore, the number of class--route incidence is

\[
 1+8+16+60+48+60=193.
\]
The incidences are not 193 distinct routes: the three neutrino flavors share the rank-twenty neutral support; the up flavors reuse a rank-sixteen support and the down flavors a rank-twenty support. The flavor projector must refine these supports. In the union of 36 quark cells, color is exactly

\[
 12_R+12_G+12_B.
\]
This equality certifies the global chromatic support, but does not yet authorize the selection of a cell by flavour.

\Needspace{7\baselineskip}
\section{Generation, Charge, Spin, and Mass-Character Decomposition of the Leptonic Sector}
\Needspace{7\baselineskip}
\subsection{Electron: center, vacancies and orientation}
The electron is exceptional because its fibre contains the unique fixed point \((5,5)\). Its reduction to scalar mass does not require breaking an orbit. The active route is

\[
 e_{\rm act}=(2,2,5,-4,-3,-2)^2,
\]
while the historical envelope

\[
 e_{\rm env}=(4,3,5,-4,-3,-5)^2
\]
conserves the sign and absolute maximum of the two leaves. The envelope retains the word

\[
 \tau_e=+++---+++---,
\]
but does not substitute active history.

Four data that must remain distinct are:

\[
 \text{raw signature}=(24,0,12,0,-12),
\]
\[
 \text{even CT--108 signature}=(-12,-4,12,-6,12),
\]
\[
 \text{event vector}=e_{\rm act},
\qquad
 \text{affine origin}=v_e=0.
\]
The basal selector \((-22,+15)\) has two complementary internal readings. The regional decomposition gives

\[
 -22=-(27-5),\qquad +15=5\cdot3=\binom62,
\]
and the framework of vacancies

\[
 V_e=
 \begin{pmatrix}
 21&22\\
 6&7
 \end{pmatrix}
\]
satisfies

\[
 \det V_e=21\cdot7-22\cdot6=15.
\]
The \(22\) measures the oriented deficit of the extension; the \(15\), the area retention that survives closure. Thus the electronic particle appears as the central persistence of an oriented vacancy, not as a mass inserted into the central cell.

The complete equation is

\[
 \mathfrak e_e^\beta=
 \frac{\sqrt3}{4}
 \left[
 \exp\!\left(\frac{\varphi}{\pi^2}\right)
 -22\alpha_{\rm HMT}^{3}
 \right]
 (1+15\Delta_4)
 R_{\rm act}^{80-54\alpha_{\rm HMT}+6\Delta_4}.
\]
In the electronic chart,

\[
 \mathfrak e_e^\beta
 =0.5109989506900008086082571648\ldots\ {\rm MeV}.
\]
\Needspace{7\baselineskip}
\subsection{Realization of the muon as a 2nd-generation leptonic section: fibre, signature and character}
The muon resides in a leptonic multisection of rank eight. Its recovered structural trace comprises the cells

\[
 \{21,23,25,39,43,57,59,61\},
\]
the counts of turn/change \(22/25\), the north--south axis,

\[
 (N,E,S,O)=(62,25,7,14),
\qquad
 (h_{369},r_5)=(57,9),
\]
and the central coordinate

\[
 (n_A,s_C)=(42,-3).
\]
The proper object is the compressed operator.

\[
 \widehat M_\mu=P_\mu\widehat M P_\mu.
\]
Its reduction to a number requires prove that the projector of flavor and generation commutes with the operator, or that the state prepared has variance null. Given the signature enriched

\[
 (42,-3,8,0,2),
\]
the evaluation of character is

\[
 105.658360294435829\ldots\ {\rm MeV}.
\]
This figure is an exact control of the evaluation of that signature; the complete physical derivation must reproduce the signature from the genealogical record selected by the coupling. The muonic width is obtained from the temporal decay operator and not from \(\widehat M_\mu\).

\Needspace{7\baselineskip}
\subsection{Realization of the tau as a 3rd-generation leptonic section: fibre, signature and character}
The tau occupies a sixteen-rank lepton multisection. Its footprint comprises

\[
 17/24,\qquad \text{east--west axis},\qquad
 (N,E,S,O)=(29,55,7,17),
\]
\[
 (h_{369},r_5)=(59,6),\qquad
 (n_A,s_C)=(64,0).
\]
The enriched signature

\[
 (64,0,25,-1,6)
\]
yields

\[
 1776.8609176314323\ldots\ {\rm MeV}.
\]
As in the muon case, the well-defined structure is the operator block; the scalar mass is obtained when the tau coupling selects a scalar block or an isolated pole.

\Needspace{7\baselineskip}
\subsection{Common neutrino fiber and flavor separation through the PMNS reader}
The three neutrino flavors share the space

\[
 \mathcal H_\nu=
 \bigoplus_{j=1}^{4}\mathcal H_{\nu,G_j},
\qquad
 \dim\mathcal H_\nu=2+4+6+8=20.
\]
The vectors

\[
 f_e,\quad f_\mu,\quad f_\tau
\]
constitute an interaction basis; they are not three mass eigenvectors. Nor is depth \(G_j\) a fourth sterile flavour: a sterile interpretation requires an additional physical projector.

The neutral extractor is defined by

\[
 T_\nu(m)=\sum_{t\in E_m}
 \omega(t)\varepsilon(t)\theta(t),
\qquad
 e_m^\nu=T_\nu(m)-T_\nu(m+6),
\]
\[
 \tau_m^\nu=\operatorname{sgn}(e_m^\nu).
\]
It is obtained from it.

\[
 b_{120}^\nu
 =\sum_{a=0}^{3}
 \operatorname{sgn}\!\left(\sum_{m\in I_a}e_m^\nu\right),
\]
\[
 b_\zeta^\nu=\sum_m\tau_m^\nu\tau_{m+3}^\nu.
\]
The neutrino signature possesses the form

\[
 v_\nu^{(i)}
 =(0,0,k_{\Delta_4}^{(i)},b_{120}^{\nu_i},b_\zeta^{\nu_i}).
\]
The inversion bidirectional changes \(b_{120}\) and preserves \(b_\zeta\); this asymmetry produces internally the two orientations of hierarchy.

\Needspace{7\baselineskip}
\subsection{Fermionic Nature of the Neutrino}
Neutrality does not determine statistics. The neutrino sector lies in the exterior completion and carries spinorial holonomy:

\[
 U_\nu(2\pi)=-I,\qquad U_\nu(4\pi)=I,
\]
\[
 \{a_{\nu,i},a_{\nu,j}^\dagger\}=\delta_{ij}.
\]
Therefore, every physical neutrino mode is fermionic. The Dirac/Majorana question is subsequent: it determines whether conjugation interchanges two distinct sections or identifies a section with its conjugate within a real representation. Neither possibility makes the neutrino a boson.

\Needspace{7\baselineskip}
\subsection{Mass operator and PMNS}
Let \(P_\nu\) be the rank-three projector onto the subspace of mass modes. First diagonalize

\[
 P_\nu\widehat M_\nu P_\nu
 =\sum_{i=1}^{3}m_i\,|n_i\rangle\langle n_i|.
\]
Then the transport between bases is constructed:

\[
 (f_e,f_\mu,f_\tau)
 =(n_1,n_2,n_3)U_{\rm PMNS}^\dagger.
\]
A full-rank kernel that reads sheet, orientation, carry, vacancy, and holonomy can be written

\[
 G_{ab}=
 \frac1{\sqrt{|L_a||N_b|}}
 \sum_{c\in L_a}\sum_{d\in N_b}K_{\rm reg}(c,d),
\]
\[
 U_{\rm PMNS}=\operatorname{polar}(G),
\qquad \det G\ne0.
\]
The internal phase of rank one produces the angles

\[
 (19.6807^\circ,56.4970^\circ,29.6224^\circ)
\]
and, for that reason, cannot be the complete physical PMNS kernel. Full rank is a mathematical requirement, not a numerical retouching.

The observables

\[
 m_1,m_2,m_3,\qquad
 \Delta m_{ij}^2,\qquad
 m_\beta,\qquad
 m_{\beta\beta},\qquad
 \sum_i m_i
\]
they are different maps of the same spectrum. They should not be collapsed into a single "neutrino mass".

\Needspace{7\baselineskip}
\section{Quarks, color, generations, and CKM mixing}
\Needspace{7\baselineskip}
\subsection{Structural position of the \(6\) quark flavors}
Each quark is defined by

\[
 \mathfrak q_X=
 (\text{charge branch},\text{generation},
 \mathbb Z_3^{\rm color},\text{multisection},
 \text{genealogical record},\text{sheet}).
\]
The structural position of the six flavors is summarized by their traces prior to any mass chart:

\Needspace{7\baselineskip}
\begingroup
\fontsize{7.2}{8.2}\selectfont
\setlength{\tabcolsep}{1.9pt}
\begin{longtable}{L{0.94cm}L{2.81cm}L{2.03cm}L{1.87cm}L{3.28cm}L{2.18cm}L{1.87cm}}
\toprule
\rowcolor{HMTNavy}\HMTtablehead{Flavor} & \HMTtablehead{Branch/\allowbreak{}generation} & \HMTtablehead{Giro/\allowbreak{}cambio} & \HMTtablehead{Eje} & \HMTtablehead{\((N,E,S,O)\)} & \HMTtablehead{\((h_{369},r_5)\)} & \HMTtablehead{\((n_A,s_C)\)} \\
\midrule
\endfirsthead
\multicolumn{7}{l}{\sffamily\scriptsize\color{HMTGray}Continuation of the table}\\[-1pt]
\toprule
\rowcolor{HMTNavy}\HMTtablehead{Sabor} & \HMTtablehead{Rama/\allowbreak{}generación} & \HMTtablehead{Giro/\allowbreak{}cambio} & \HMTtablehead{Eje} & \HMTtablehead{\((N,E,S,O)\)} & \HMTtablehead{\((h_{369},r_5)\)} & \HMTtablehead{\((n_A,s_C)\)} \\
\midrule
\endhead
\midrule
\endfoot
\bottomrule
\endlastfoot
\(u\) & positiva/\allowbreak{}1 & \(21/\allowbreak{}30\) & N--S & \((34,31,31,12)\) & \((51,13)\) & \((11,7)\) \\
\rowcolor{HMTLight}\(d\) & negativa/\allowbreak{}1 & \(20/\allowbreak{}32\) & E--O & \((27,35,18,28)\) & \((47,17)\) & \((17,8)\) \\
\(c\) & positiva/\allowbreak{}2 & \(21/\allowbreak{}29\) & E--O & \((20,38,24,26)\) & \((49,16)\) & \((61,8)\) \\
\rowcolor{HMTLight}\(s\) & negativa/\allowbreak{}2 & \(17/\allowbreak{}27\) & N--S & \((50,32,12,14)\) & \((57,11)\) & \((41,-3)\) \\
\(t\) & positiva/\allowbreak{}3 & \(21/\allowbreak{}26\) & N--S & \((46,27,16,19)\) & \((57,9)\) & \((100,-1)\) \\
\rowcolor{HMTLight}\(b\) & negativa/\allowbreak{}3 & \(21/\allowbreak{}30\) & E--O & \((12,63,9,24)\) & \((47,7)\) & \((71,-5)\) \\
\end{longtable}
\endgroup
Estas posiciones no son seis números elegidos: son seis clases de representación que actúan sobre tripletes de color. La ruta física puede ser una órbita o un bloque de rutas; no tiene que ser un punto.

\Needspace{7\baselineskip}
\subsection{Invarianza escalar del carácter de masa bajo la acción de color}
Sea \(P_q\) el proyector de sabor y generación y \(P_{\rm RGB}\) el proyector a la órbita de color. La masa quark es

\[
 \widehat M_q=
 P_qP_{\rm RGB}\widehat M P_{\rm RGB}P_q.
\]
If

\[
 [\widehat M_q,U(g)]=0,\qquad g\in SU(3),
\]
then

\[
 \widehat M_q=m_qI_3
\]
over the triplet irreducible. This equality explains by which the flavor carries a mass common although their three colors occupy positions distintas in the atlas.

\Needspace{7\baselineskip}
\subsection{Chart of schema and scale}
The internal HMT coordinate must be transported to a renormalized definition:

\[
 m_q^{\mathcal S}(\mu)
 =Z_m^{\mathcal S}(\mu,\mu_0)\,
 \mathcal R_q(\widehat M_q),
\]
with composition law

\[
 Z_m(\mu_2,\mu_0)
 =Z_m(\mu_2,\mu_1)Z_m(\mu_1,\mu_0).
\]
For \(u,d,s\), the natural comparison uses \(\overline{\rm MS}\) at \(2\,{\rm GeV}\); for \(c,b\), the scale of the renormalized mass itself; for \(t\), one must state whether the chart is a pole, MSR, or kinematic chart. Without this transport, a numerical difference is not yet a residual of the law.

\Needspace{7\baselineskip}
\subsection{Assessments of the declared signatures}
The historical central coordinates produce:

\Needspace{7\baselineskip}
\begingroup
\footnotesize
\setlength{\tabcolsep}{2.1pt}
\begin{longtable}{L{4.68cm}L{10.61cm}}
\toprule
\rowcolor{HMTNavy}\HMTtablehead{Flavor} & \HMTtablehead{Evaluation} \\
\midrule
\endfirsthead
\multicolumn{2}{l}{\sffamily\scriptsize\color{HMTGray}Continuation of the table}\\[-1pt]
\toprule
\rowcolor{HMTNavy}\HMTtablehead{Sabor} & \HMTtablehead{Evaluación} \\
\midrule
\endhead
\midrule
\endfoot
\bottomrule
\endlastfoot
\(u\) & \(2.165924\allowbreak{}536173\allowbreak{}907\ldots\ {\rm MeV}\) \\
\rowcolor{HMTLight}\(d\) & \(4.679550\allowbreak{}162948\allowbreak{}874\ldots\ {\rm MeV}\) \\
\(c\) & \(1.270500\allowbreak{}838641\allowbreak{}911\ldots\ {\rm GeV}\) \\
\rowcolor{HMTLight}\(s\) & \(92.940093\allowbreak{}907845\allowbreak{}64\ldots\ {\rm MeV}\) \\
\(t\) & \(172.597479\allowbreak{}544019\allowbreak{}1\ldots\ {\rm GeV}\) \\
\rowcolor{HMTLight}\(b\) & \(4.190119\allowbreak{}763299\allowbreak{}790\ldots\ {\rm GeV}\) \\
\end{longtable}
\endgroup
Estas evaluaciones comprueban el carácter una vez dada la firma. La deducción completa exige obtener cada firma, su refinamiento de torre y su carta de esquema a partir del proyector de sabor y del registro genealógico correspondiente.

\Needspace{7\baselineskip}
\subsection{Determinación de los ángulos CKM desde la reflexión racional de las rutas quark}
Sean

\[
 A=1000\alpha_{\rm HMT},\qquad
 h=\frac{C^*}{6},\qquad
 \Delta=\frac{180}{\pi}\Delta_4.
\]
The angular grammar gives

\[
 \begin{pmatrix}
 \theta_{12}\\
 \theta_{23}\\
 \theta_{13}
 \end{pmatrix}
 =
 \begin{pmatrix}
 2&-5&28\\
 0&7&-25/2\\
 1&-20&-2/3
 \end{pmatrix}
 \begin{pmatrix}
 A\\h\\\Delta
 \end{pmatrix},
\qquad
 \delta_{\rm CP}=9A.
\]
The exact linear falsifier is

\[
 3857\,\delta_{\rm CP}
 =9\bigl(
 1528\theta_{12}
 +3380\theta_{23}
 +801\theta_{13}
 \bigr).
\]
From the three angles and \(\delta_{\rm CP}\), the unitary matrix \(V_{\rm CKM}\) is constructed. Its function in the mass law is not to produce eigenvalues, but to transport the down operators to the up basis:

\[
 B_d^{(u)}=V_{\rm CKM}B_dV_{\rm CKM}^*.
\]
\Needspace{7\baselineskip}
\subsection{Jarlskog Invariant}
The CP violation invariant is

\[
 J_{\rm CKM}
 =c_{12}c_{23}c_{13}^2
 s_{12}s_{23}s_{13}\sin\delta_{\rm CP},
\]
and the internal evaluation is

\[
 J_{\rm CKM}
 =3.1189723326632947\times10^{-5}.
\]
Equivalently, the compatibility between spectra and mixing may is read in the determinant of the commutator of the operators Hermitian up/down. A value not null of \(J_{\rm CKM}\) records orientation physical not removable; is not a corrector of mass added to each quark.

\Needspace{7\baselineskip}
\section{Gauge and scalar sectors}
\Needspace{7\baselineskip}
\subsection{Chart of the null sector: mass routes \(0\), gauge representations, and persistence}
Let

\[
 \mathcal A_{\rm sp}=\mathbb R[R]/(R^2-1),
 \qquad
 P_\pm=\frac{1\pm R}{2}.
\]
For

\[
 x=r_+P_++r_-P_-,
\]
the split norm is

\[
 N_{\rm sp}(x)=r_+r_-.
\]
The null section is a direction of zero norm. Its mass is exactly

\[
 m=0,
\]
without passing through a limit of positive characters and without assigning \(\kappa=0\). Nullity belongs to the type of the section.

\Needspace{7\baselineskip}
\subsection{Realization of the photon in the electromagnetic null sector: helicity, gauge, and character}
The photon is the transverse vector realization of a null gauge section. The physical projector must leave exactly two polarizations:

\[
 \operatorname{rg}P_\gamma^{\rm trans}=2,
\qquad
 P_\gamma^{\rm trans}k=0.
\]
The rest mass is

\[
 m_\gamma=0.
\]
Bidirectional propagation may carry an oriented anisotropy.

\[
 \epsilon=q_-^{90}-q_+^{90},
\qquad
 c_\pm=\frac{c}{1\pm\epsilon},
\]
while the harmonic mean preserves

\[
 \frac{2}{1/c_++1/c_-}=c.
\]
Thus, the outward and return paths may differ without altering the two--way invariant. This bias is not a photon mass.

A distinct Proca/Stückelberg realization appears if a spectral gap in the genealogical record satisfies

\[
 \omega^2=k^2+m_T^2
\]
and there exists a dimensional entangler

\[
 m_T\longmapsto m_\gamma.
\]
The induced delay would obey

\[
 \frac{\Delta T}{T}
 \simeq\frac12
 \left(\frac{m_\gamma c^2}{\hbar\omega}\right)^2.
\]
The null Maxwell sector and the massive Proca realization are distinct branches; neither invalidates the other.

\Needspace{7\baselineskip}
\subsection{Chromodynamic realization of the null sector: gluon octet and adjoint representation of \(SU(3)\)}
The residual trit generates the color space, whose traceless part realizes

\[
 \mathfrak{sl}_3\longrightarrow\mathfrak{su}_3.
\]
The gluon occupies the eight-dimensional adjoint representation. As long as color symmetry remains exact,

\[
 m_g=0.
\]
Confinement, screening, and the collective spectral gap are properties of the interacting spectrum; they do not constitute an elementary gluon pole mass. An effective mass must specify the observable, gauge, regime, and operator that defines it.

\Needspace{7\baselineskip}
\subsection{Realization of the bosons \(W^\pm\): charge, polarization, and electroweak character}
The sector \(W\) is a charged vector block with conjugation

\[
 C_M:W^+\longleftrightarrow W^-.
\]
Its historical footprint is

\[
 15/37,\qquad \text{east--west axis},\qquad
 (n_A,s_C)=(94,-1),
\]
and the enriched signature declared is

\[
 (94,-1,0,-2,4).
\]
The corresponding evaluation is

\[
 80.37896490970455\ldots\ {\rm GeV}.
\]
The realization complete is the pole of the resolvente of the block \(P_W\widehat M P_W\); its part imaginaria determines the width. The degeneration of \(W^+\) and \(W^-\) proceeds of the conjugation of the operator.

\Needspace{7\baselineskip}
\subsection{Realization of the \(Z^0\) boson: neutral mixing, polarization, and electroweak character}
The \(Z^0\) is the neutral vector block of the electroweak coupling. Its historical trace is

\[
 20/23,\qquad \text{east--west axis},\qquad
 (n_A,s_C)=(95,-1),
\]
and the enriched signature

\[
 (95,-1,-11,0,-4).
\]
The evaluation of the character gives

\[
 91.18753344431341\ldots\ {\rm GeV}.
\]
The neutral projector must be constructed before the pole chart; mass and width are two coordinates of the same pole, not a single figure.

\Needspace{7\baselineskip}
\subsection{Realization of the Higgs sector: scalar mode, coupling and mass character}
The Higgs is a persistent scalar section of the gauge block with declared CP parity. Its trace is

\[
 24/29,\qquad \text{north--south axis},\qquad
 (n_A,s_C)=(97,9).
\]
The conditional evaluation of the character produces

\[
 125.2924660425361\ldots\ {\rm GeV}.
\]
The HMT definition does not identify “scalar” with “mass”: zero spin fixes the rotational representation; mass follows from the spectrum of the coupled scalar block. The width and decay channels require the temporal operator.

\Needspace{7\baselineskip}
\subsection{The anomalous magnetic moment of the muon}
The internal functional

\[
 a_\mu=
 \frac{\alpha_{\rm HMT}}{2\pi}
 +\frac{\alpha_{\rm HMT}(\varphi-1)}{1000}
 +\frac{\alpha_{\rm HMT}^2}{1000(54+1)}
\]
gives

\[
 a_\mu=0.001165920713008014\ldots,
\]
and

\[
 g_\mu=2(1+a_\mu)
 =2.002331841426016028\ldots.
\]
The three terms correspond, respectively, to the basal rotation, the oriented return in the completion \(1000\), and the quadratic correction mediated by the defect \(54\). This observable verifies the internal structure of the muonic sector but does not by itself select its mass route.

\Needspace{7\baselineskip}
\section{Construction of composite particles through non-additive composition of paths and bonds}
\Needspace{7\baselineskip}
\subsection{Composition operator}
Route composition is not strictly additive because boundary, temporal order, and carry interact. For two ledgers,

\[
 \mathcal L_1\star\mathcal L_2
 =\mathcal L_1+\mathcal L_2
 +\mathfrak b(\mathcal L_1,\mathcal L_2),
\]
where \(\mathfrak b\) is the binding cocycle. Associativity requires

\[
 \mathfrak b(L_1,L_2)
 +\mathfrak b(L_1\star L_2,L_3)
 =
 \mathfrak b(L_2,L_3)
 +\mathfrak b(L_1,L_2\star L_3).
\]
The mass reader acts after this composition:

\[
 (\Gamma_i,\mathcal T,\partial)
 \xrightarrow{\mathfrak C_{12}}
 \mathcal L_{\rm comp}
 \xrightarrow{L_{\rm tick}}
 \sigma_{\rm comp}
 \xrightarrow{\chi_{\rm full}}
 \widehat M_{\rm comp}.
\]
\Needspace{7\baselineskip}
\subsection{Mesonic Realization via Quark–Antiquark Composition and Cocylic Link Law}
The mesonic domain contains 432 oriented expressions.

\[
 M(q,q')=q\otimes(q')^\vee
\]
with color compatible. The physical projector imposes color singlet and fixes

\[
 (J^{PC},I,I_3,S,C,B,\ldots),
\]
together with the temporal linking tree. The mass of a meson is not \(m_q+m_{\bar q}\); it is an eigenvalue of the composite ledger.

The six historical control evaluations include

\[
 m_{\pi^0}=134.9767243278721\ldots\ {\rm MeV},
\]
\[
 m_{\pi^\pm}=139.5704851715722\ldots\ {\rm MeV},
\]
\[
 m_{K^\pm}=493.6770856495306\ldots\ {\rm MeV},
\]
\[
 m_{K^0}=497.6111864977505\ldots\ {\rm MeV}.
\]
These figures are exact evaluations of declared records. Genealogical closure state by state requires regenerating each record by means of \(\mathfrak C_{12}\), including its binding boundary.

\Needspace{7\baselineskip}
\subsection{Baryonic realization via trilinear composition of quark paths and color neutralization}
The baryonic domain contains 1728 RGB expressions. The physical projector must simultaneously impose:

\begin{enumerate}
\item color singlet;
\item fermionic statistics;
\item flavor and isospin;
\item spin and symmetrization;
\item triple tree with carry and memory.
\end{enumerate}
For proton and neutron, the historical signatures

\[
 \sigma_p=(59,0,9,1,0),\qquad
 \sigma_n=(59,1,-34,0,1)
\]
give

\[
 m_p=938.2717515469849\ldots\ {\rm MeV},
\]
\[
 m_n=939.5658104725070\ldots\ {\rm MeV}.
\]
Apparent precision does not replace derivation of the genealogical binding
record. The sign and magnitude of

\[
 m_n-m_p
\]
must follow from orientation, isospin, vacancies, and the linking boundary before external comparison is opened.

\Needspace{7\baselineskip}
\subsection{Excited states and hadronic resonances}
A composite identity may possess several poles. The physical state is individualized by

\[
 (\text{constituents},\mathcal T,\partial,
 J^{PC},I,\text{radial},\text{orbital},\text{channel}).
\]
The same flavor composition is not sufficient to distinguish excitations. Each pole

\[
 z_k=M_k-\frac{i}{2}\Gamma_k
\]
must proceed from an eigenvector of the composite operator and from a decay channel. This extends the law from a few stable masses to the entire spectrum of mesons and baryons without converting the external catalog into a generator.

\Needspace{7\baselineskip}
\subsection{Comprehensive inventory of widths and transition rates of the atlas realizations}
Let \(P\) be the prepared subspace and \(Q=I-P\). The effective Hamiltonian is

\[
 H_{\rm eff}(z)
 =PHP+PHQ(z-QHQ)^{-1}QHP.
\]
The poles of

\[
 \det(z-H_{\rm eff}(z))=0
\]
determine masses and widths. The mass character supplies the basal real part of the operator; the self-energy term of \(Q\) supplies shift and decay. Therefore, the 384 widths of the inventory require a channel reader in addition to the reader of the 471 masses.

\Needspace{7\baselineskip}
\section{Topological sectors and hypothetical particles}
\Needspace{7\baselineskip}
\subsection{Fibonacci topological sector: ring-based, fusion, and braid representation}
The areal--volumetric incidence

\[
 F_{\rm av}=
 \begin{pmatrix}
 0&1\\1&1
 \end{pmatrix}
\]
realize the ring

\[
 \tau^2=1+\tau,\qquad
 \operatorname{FPdim}\tau=\varphi.
\]
The primary output is a fusion rule and a quantum dimension. A Fibonacci anyonic particle also requires \(F\) and \(R\) matrices, the pentagon and hexagon equations, chirality, a spectrally gapped Hamiltonian, and a map from the TPK state to physical space. Only then can one define

\[
 m_{\rm brecha espectral}=\frac{\Delta E}{c_{\rm eff}^2}.
\]
\Needspace{7\baselineskip}
\subsection{Ising–Majorana Sector: Simple Objects, Fusion Rules, and Braid Representation}
The Ising category satisfies

\[
 \psi^2=1,\qquad
 \psi\sigma=\sigma,\qquad
 \sigma^2=1+\psi,\qquad
 d_\sigma=\sqrt2.
\]
Four levels must be distinguished:

\begin{enumerate}
\item the route involution \(C_M=-\operatorname{rev}\);
\item the parity of the CAR completion;
\item a self-conjugate relativistic spinor;
\item a topological Ising zero mode.
\end{enumerate}
A relativistic Majorana fermion requires a real spinor representation, a dynamics, and a mass term. A Majorana zero mode is characterized by a spectral gap, splitting, and protection; it has no universal mass deducible solely from the fusion ring. The identity of a route with its conjugate is a necessary, but not sufficient, condition for a Majorana realization.

\Needspace{7\baselineskip}
\subsection{Braid representation on \(\mathbb Z/6\mathbb Z\) and topological characters}
Let

\[
 w_6:B_n\longrightarrow\mathbb Z/6,
\qquad
 \chi_k(b)=e^{2\pi ikw(b)/6}.
\]
The representations

\[
 \rho_q(\sigma_i)=qP_i
\]
descend to the symmetric group only when \(q^2=1\). The other four values are genuinely braided. Their mathematical existence does not imply four particles: a physical realization requires a Hamiltonian, a spectral gap, braid operators, stability, and an intertwiner with a multisection.

\Needspace{7\baselineskip}
\subsection{Impossibility of a Persistent Tachyon}
The action line is positive,

\[
 \hbar_{\rm int}>0,
\]
the character satisfies

\[
 \chi_{\rm full}>0,
\]
and the beta operator is obtained by positive conjugation:

\[
 \widehat M^\beta
 =R_{\rm act}^{\widehat K/2}
 \widehat M^{(0)}
 R_{\rm act}^{\widehat K/2}\ge0.
\]
Therefore,

\[
 \operatorname{Spec}(\widehat M^\beta)\subset[0,\infty).
\]
A persistent stable section cannot generate

\[
 m^2<0.
\]
If the linearization of a state produces negative curvature or an imaginary frequency, that state is an instability that leaves the section, not a superluminal particle. If the prolongation terminates, it is a virtual section with a terminal obstruction. Thus, “tachyon” decomposes into two well-typed HMT objects ---instability or virtuality--- and does not constitute a persistent species.

\Needspace{7\baselineskip}
\subsection{Gravity without a mandatory elementary spin \(2\) particle}
The gravitational realization follows the chain

\[
 \text{route holonomy}
 \longrightarrow(e,\omega)
 \longrightarrow(T,R)
 \longrightarrow
 S_{\rm PCH},
\]
where \(e\) is the costructure, \(\omega\) the connection,

\[
 T=de+\omega\wedge e,\qquad
 R=d\omega+\omega\wedge\omega,
\]
and \(S_{\rm PCH}\) is the Palatini--Cartan--Holst action. Gravity appears as the global torsion-and-curvature response of the persistent network. The collective vibration of the model universe may have a spin-two tensor component, but HMT need not posit an additional elementary particle to carry gravity.

The internal gravitational constant is typed through

\[
 G_{\rm HMT}
 =\frac{c_{\rm int}^3}{\hbar_{\rm int}}
 \bigl(\delta_\theta\Lambda_G\bigr)^2,
\]
where \(\delta_\theta\) is the elementary angular torsion and \(\Lambda_G\) is the spectral length selected by the gravitational holonomy. This equation converts torsion and route scale into gravitational coupling without using a mass as normalization. The total vibration is represented by the Hessian

\[
 \mathcal H_{\rm univ}
 =\operatorname{Hess}_{(e,\omega)}S_{\rm PCH}
\]
over the global state. Its modes are collective excitations of the model universe; a tensorial component is not automatically a separate particle.

An excitation that is to be called a "graviton" must demonstrate three things:

\begin{enumerate}
\item a persistent tensorial subspace;
\item an isolated pole of the linearized operator;
\item a physical coupling that does not duplicate the torsion already present.
\end{enumerate}
Without those three conditions, the correct object is the collective mode of the geometry, not a new mass row.

\Needspace{7\baselineskip}
\subsection{Barbero-Immirzi and the hexad-octad link}
With

\[
 A=1000\alpha_{\rm HMT},
\qquad
 q_\pm=\exp\!\left[-\frac{\pi}{180}(A\pm C^*)\right],
\]
the hexadic-octadic functional uses two distinct incidences of the temporal towers. The combination

\[
 \varepsilon_{\rm BI}
 =12\Delta S_{90}-\Delta S_{120}
\]
and the angular term produce

\[
 \gamma_{\rm BI}
 =\frac{\pi}{180}AC^*
 +\varepsilon_{\rm BI}
 =0.2740076726944593\ldots.
\]
The weight \(12:-1\) comes from reading the twelve dodecaphasic sectors
and the unique oriented return seam. The equation
\(4a-3b=45\) subsequently verifies that the pair \((12,1)\) is the minimal
positive integer solution. This object must not be confused with the dodecaphasic mass extractor or with the angular residue \(90\)--\(120\); it shares the harmonic operator, but has a different domain and function.

\Needspace{7\baselineskip}
\subsection{Extensions of the atlas toward dark sectors: domains, readers and physical status}
A dark sector HMT is a persistent multisection with

\[
 P_{\rm vis}P_{\rm dark}=0
\]
for the acoplamientos electromagnetico and cromatico, but with posible incidencia gravitatoria or of boundary:

\[
 P_{\rm grav}P_{\rm dark}\ne0.
\]
Its “darkness” is a property of the coupling kernel, not a mass label. The law first predicts an operator and its spectrum; only a physical channel determines whether the mode is stable, metastable, or resonant. This definition makes it possible to search for invisible multisections without inventing a mass from the observed value of a cosmological density.

\Needspace{7\baselineskip}
\subsection{Sterile neutrino section: compatibility of fiber, mixing, and absence of gauge charge}
A sterile neutrino is not the depth \(G_4\). It is an eigenvector \(n_s\in\mathcal H_\nu\) satisfying

\[
 P_{\rm weak}n_s=0,\qquad
 P_{\rm mass}n_s=n_s.
\]
It may mix with the three active modes through an enlarged neutral kernel, but its existence requires

\[
 \dim\ker(P_{\rm weak}|_{\mathcal H_\nu})>0
\]
and an eigenvalue of mass in ese kernel. This definition produces a criterion of search internal: diagonalize first the block neutral and check afterward its coupling weak.

\Needspace{7\baselineskip}
\subsection{Axionic and pseudoscalar sections: topological character and edge coupling}
A pseudoscalar axion candidate is a section whose orientation changes under CP,

\[
 CP\,a\,CP^{-1}=-a,
\]
and whose displacement compensates a phase topological of the sector of color. In HMT must correspond to a way of torsion/holonomy odd and to a projector that intertwines ese way with the density topological of gauge. Its mass is obtained of the second derivative of the potential effective,

\[
 m_a^2=
 \left.\frac{\partial^2V_{\rm eff}}{\partial a^2}\right|_{a=a_0},
\]
once \(V_{\rm eff}\) has been constructed from the genealogical record of vacancies and boundary. The mere existence of a channel \(C^*\mapsto-C^*\) is not enough to assert an axionic particle.

\Needspace{7\baselineskip}
\subsection{Dark Photon and Hidden Massive Vector}
A dark photon requires a second gauge section \(A'\) and a mixing coupling

\[
 \mathcal L_{\rm mix}
 =-\frac{\varepsilon_{\rm kin}}2F_{\mu\nu}F'^{\mu\nu}.
\]
In terms of fibers, \(\varepsilon_{\rm kin}\) is an intertwiner between two null subspaces. If the second sector acquires a Stueckelberg or boundary spectral gap,

\[
 \widehat M_{A'}^2\ne0,
\]
a massive dark vector appears. HMT predicts that the mixing parameter must follow from a common memory incidence; it cannot be fixed by the mass it is intended to explain.

\Needspace{7\baselineskip}
\subsection{WIMP and stable dark state}
A WIMP is a dark sector that also satisfies:

\[
 [P_{\rm weak},P_X]\ne0,\qquad
 [P_{\rm em},P_X]=[P_{\rm color},P_X]=0,
\]
and possesses a stable mode

\[
 |\lambda_X|=1.
\]
Stability may be protected by a parity of the genealogical record. Its mass is an eigenvalue of the weakly coupled dark block. “WIMP” does not constitute a new mathematical family: it is a joint condition on coupling and persistence over a multisection.

\Needspace{7\baselineskip}
\subsection{Monopole Section: Magnetic Charge, Dirac Condition, and Topological Compatibility}
A monopole is defined by a non-trivial topological class of the gauge bundle:

\[
 \frac1{2\pi}\int_{S^2}F=n\ne0.
\]
HMT must realize this integer as the net holonomy of a closed family of routes. Its mass, if the mode is dynamical and stable, belongs to the spectrum of the topological-defect operator. A point vacancy or a nonzero carry is not by itself a monopole; the descent to the integral class must be proved.

\Needspace{7\baselineskip}
\subsection{Supersymmetric Extensions of the Atlas: Representation Matching and Characters}
An HMT supersymmetry would require an odd operator \(Q\) that exchanges the exterior and symmetric completions:

\[
 Q:\mathcal F_-\leftrightarrow\mathcal F_+,
\qquad
 Q^2=H,
\]
and that interweaves the law of masses,

\[
 Q\widehat M_-=\widehat M_+Q.
\]
Only then would the bosonic and fermionic spectra be paired. The separate existence of CAR and CCR does not establish supersymmetry or by itself predict selectrons, neutralinos, or gluinos. Each partner requires a nonzero intertwiner and a symmetry-breaking rule that determines its mass splitting.

\Needspace{7\baselineskip}
\subsection{Tetraquarks, pentaquarks, glueballs, and hybrids}
The composition grammar is extended without changing the law:

\[
 \mathfrak C_{12}^{(n)}:
 \Gamma_1\times\cdots\times\Gamma_n
 \times\mathcal T_n\times\partial_n
 \longrightarrow\mathcal L_{\rm comp}.
\]
A tetraquark uses \(qq\bar q\bar q\); a pentaquark, \(qqqq\bar q\); a glueball, records of the color adjoint projected onto the singlet; a hybrid, quark and gauge routes in the same tree. The physical classification depends on the color, spin, parity, and linkage-tree projector, not merely on the list of constituents.

\Needspace{7\baselineskip}
\subsection{Dilaton and radion sectors: volume and compactification scalar modes}
The dilaton is a way of the section of scale; the radion, a way that modifies a length or volume internal. In the theory of mass appear as fluctuations of the chart and of the memory geometric:

\[
 t_0\mapsto e^{\phi_d}t_0,\qquad
 \mathcal V_{\rm int}\mapsto e^{\phi_r}\mathcal V_{\rm int}.
\]
To become particles they must possess a kinetic term, a potential, a persistent projector, and a pole. So long as they do not possess them, they are collective degrees of the dimensional chart, not additional scalar species.

\Needspace{7\baselineskip}
\subsection{Exclusion of Physical Ghosts}
The spectral measure of a physical state is positive:

\[
 \mu_{P,a}^{\rm mass}(B)\ge0.
\]
A negative-norm direction does not belong to the prepared state space. If a parametrization introduces a “ghost,” the physical realization must quotient out the gauge redundancy or reject the coupling. This exclusion is independent of the tachyonic prohibition: a ghost violates positivity of the norm; a persistent tachyon would violate positivity of the mass operator.

\Needspace{7\baselineskip}
\section{Universal Inventory of Masses and Widths}
\Needspace{7\baselineskip}
\subsection{Posterior correspondence between the HMT–MD realizations and the external physical census}
The contrast inventory contains

\[
 1170\ \text{identities},\qquad
 438\ \text{identity groups},
\]
\[
 471\ \text{mass observables},\qquad
 384\ \text{width observables}.
\]
Its partition is

\Needspace{7\baselineskip}
\begingroup
\footnotesize
\setlength{\tabcolsep}{2.1pt}
\begin{longtable}{L{5.04cm}L{5.04cm}L{5.04cm}}
\toprule
\rowcolor{HMTNavy}\HMTtablehead{Domain} & \HMTtablehead{Masses} & \HMTtablehead{Widths} \\
\midrule
\endfirsthead
\multicolumn{3}{l}{\sffamily\scriptsize\color{HMTGray}Table continued}\\[-1pt]
\toprule
\rowcolor{HMTNavy}\HMTtablehead{Dominio} & \HMTtablehead{Masas} & \HMTtablehead{Anchuras} \\
\midrule
\endhead
\midrule
\endfoot
\bottomrule
\endlastfoot
partículas generales y leptónicas & 53 & 9 \\
\rowcolor{HMTLight}quarks & 8 & 1 \\
mesones & 243 & 225 \\
\rowcolor{HMTLight}bariones & 166 & 149 \\
registro gravitatorio de contraste & 1 & 0 \\
\rowcolor{HMTLight}\textbf{Total} & \textbf{471} & \textbf{384} \\
\end{longtable}
\endgroup
Las 471 masas no son 471 especies elementales: incluyen estados excitados, resonancias, distintas cargas y múltiples realizaciones compuestas. Las 384 anchuras son observables temporales y no deben ser leídas por el carácter de masa.

\Needspace{7\baselineskip}
\subsection{Identidad interna \(\mathcal I_X\) de una realización material y criterio de individualización frente a observaciones múltiples}
Cada estado comparado debe poseer una identidad interna

\[
 \mathcal I_X=
 (\mathcal C_X,\mathcal R_X,\Gamma_X,\mathcal L_X,
 \sigma_X,\eta_X,P_X,\mathfrak M_X,\mathcal T_X),
\]
where:

\begin{itemize}
\item \(\mathcal C_X\) is the physical type of realization determined by object, representation, and codomain;
\item \(\mathcal R_X\) is the physical representation;
\item \(\Gamma_X\) is an oriented route, orbit, or composition;
\item \(\mathcal L_X\) is the genealogical record;
\item \(\sigma_X\) is the complete signature;
\item \(\eta_X\) is the tower;
\item \(P_X\) is the coupling projector;
\item \(\mathfrak M_X\) is the measurement channel;
\item \(\mathcal T_X\) is the composition tree when applicable.
\end{itemize}
Two external observations belong to the same HMT state when they share this identity and differ only in experiment, chart, or estimator. The existence of two measurements does not create two particles.

\Needspace{7\baselineskip}
\subsection{Comparison morphism between internal invariants and metrological observables}
The comparison requires the diagram

\[
 \mathcal L_M^{\rm HMT}
 \xrightarrow{\;\varsigma_{\mathcal S,\mu}\;}
 \mathbb R_{>0}u_M^{\mathcal S,\mu}
 \xleftarrow{\;\text{measurement}\;}
 \mathcal O_X^{\rm ext}.
\]
A scientifically complete row contains:

\[
 (m_X^{\rm HMT},m_X^{\rm ext},
 u,\mathcal S,\mu,\Sigma_{\rm exp},\Sigma_{\rm HMT},
 \text{pole convention}).
\]
For quarks are mandatory scheme and scale; for resonances, the convention of pole; for limits, the level of confidence; for observables correlated, the matrix of covariance.

\Needspace{7\baselineskip}
\subsection{Explicit evaluation of \(17\) representative realizations of the atlas}
The following evaluations preserve the separation between law and signature data:

\Needspace{7\baselineskip}
\begingroup
\footnotesize
\setlength{\tabcolsep}{2.1pt}
\begin{longtable}{L{4.68cm}L{10.61cm}}
\toprule
\rowcolor{HMTNavy}\HMTtablehead{Status} & \HMTtablehead{Value of the character in its chart} \\
\midrule
\endfirsthead
\multicolumn{2}{l}{\sffamily\scriptsize\color{HMTGray}Continuation of the table}\\[-1pt]
\toprule
\rowcolor{HMTNavy}\HMTtablehead{Status} & \HMTtablehead{Value of the character in its chart} \\
\midrule
\endhead
\midrule
\endfoot
\bottomrule
\endlastfoot
\(\mu\) & \(105.658360\allowbreak{}294435\allowbreak{}829\ldots\ {\rm MeV}\) \\
\rowcolor{HMTLight}\(\tau\) & \(1776.860917\allowbreak{}631432\allowbreak{}3\ldots\ {\rm MeV}\) \\
\(u\) & \(2.165924\allowbreak{}536173\allowbreak{}907\ldots\ {\rm MeV}\) \\
\rowcolor{HMTLight}\(d\) & \(4.679550\allowbreak{}162948\allowbreak{}874\ldots\ {\rm MeV}\) \\
\(c\) & \(1.270500\allowbreak{}838641\allowbreak{}911\ldots\ {\rm GeV}\) \\
\rowcolor{HMTLight}\(s\) & \(92.940093\allowbreak{}907845\allowbreak{}64\ldots\ {\rm MeV}\) \\
\(t\) & \(172.597479\allowbreak{}544019\allowbreak{}1\ldots\ {\rm GeV}\) \\
\rowcolor{HMTLight}\(b\) & \(4.190119\allowbreak{}763299\allowbreak{}790\ldots\ {\rm GeV}\) \\
\(W\) & \(80.378964\allowbreak{}909704\allowbreak{}55\ldots\ {\rm GeV}\) \\
\rowcolor{HMTLight}\(Z\) & \(91.187533\allowbreak{}444313\allowbreak{}41\ldots\ {\rm GeV}\) \\
\(H\) & \(125.292466\allowbreak{}042536\allowbreak{}1\ldots\ {\rm GeV}\) \\
\rowcolor{HMTLight}\(\pi^0\) & \(134.976724\allowbreak{}327872\allowbreak{}1\ldots\ {\rm MeV}\) \\
\(\pi^\pm\) & \(139.570485\allowbreak{}171572\allowbreak{}2\ldots\ {\rm MeV}\) \\
\rowcolor{HMTLight}\(K^\pm\) & \(493.677085\allowbreak{}649530\allowbreak{}6\ldots\ {\rm MeV}\) \\
\(K^0\) & \(497.611186\allowbreak{}497750\allowbreak{}5\ldots\ {\rm MeV}\) \\
\rowcolor{HMTLight}\(p\) & \(938.271751\allowbreak{}546984\allowbreak{}9\ldots\ {\rm MeV}\) \\
\(n\) & \(939.565810\allowbreak{}472507\allowbreak{}0\ldots\ {\rm MeV}\) \\
\end{longtable}
\endgroup
The beta electron does not belong to this table because its central route, basal coefficients, and complete reader are determined by the construction itself.

\Needspace{7\baselineskip}
